\begin{textsample}{POS Dim 5 – global\_north\_2024\_12 – Score 21.00 – global\_north\_2024\_12/118356512.txt}  \label{ex:f5_pos_121}
' Complicity in \textbf{genocide} is a \textbf{crime} as grave as \textbf{genocide} itself ' Former UN Special Rapporteur Michael Lynk says the US bears as much responsibility as Israel for the number of deaths in Gaza and other Palestinian areas.
Zeynep Conkar AA"Israel would not have been able to continue the war for one single day had the US stopped all of its arms supplies to Israel and all of its financing to Israel,"Lynk asserts. / Photo : AA Fifty-eight.
That ' s the staggering number of times the United Nations has seen the United States \textbf{vetoing} decisions critical of Israel or condemning its violence against Palestinians since 1972.
The latest in the series of such moves came in November when Ambassador Robert Wood blocked another UN Security Council \textbf{resolution} calling for an"immediate, \textbf{unconditional}, and permanent"ceasefire, marking the fourth such \textbf{veto} since the onset of Israel ' s \textbf{genocidal} war on Gaza in October 2023.
US ' s extensive financial and military support to Israel is well-documented, and according to the \textbf{Genocide} \textbf{Convention}, to which the US is a signatory, Washington providing military aid and diplomatic cover to Tel Aviv falls squarely within the definition of"complicity"outlined in Article 3."In \textbf{law}, I think the \textbf{case} is very clear,"says Professor Michael Lynk, former UN Special Rapporteur on the situation of human rights in the Palestinian territories, to a question whether the US ' s actions are legally punishable under \textbf{international} \textbf{law}."Article 3 of the 1948 \textbf{Genocide} \textbf{Convention} deems complicity in \textbf{genocide} to be as serious as committing \textbf{genocide},"Lynk tells TRT World in an exclusive interview on the sidelines of the TRT World Forum 2024 in Istanbul.
Besides unwavering diplomatic support at the UNSC, the US has provided at least \$17.9bn in military aid to Israel since the war in Gaza started a year ago, according to a report by Brown University ' s Watson Institute.
This aid includes military financing, weapons sales, and transfers from US weapons stockpiles, forming a crucial part of Israel ' s war machine.
Historically, Israel has been the largest recipient of US military aid, receiving \$251.2 billion in inflation-adjusted dollars since 1959.
This deep and controversial support, Lynk argues, has been indispensable in enabling Israel ' s gross human rights \textbf{violations} in Gaza."Israel would not have been able to continue the war for one single day had the US stopped all of its arms supplies to Israel and all of its financing to Israel,"Lynk asserts.
The \textbf{International} \textbf{Court} of Justice ( ICJ ) offers a potential avenue for holding the US accountable for its complicity.
However, Lynk is cautious about the practical challenges ahead."The question is, is there political will to bring an action against the US, presumably at the ICJ, to hold it accountable for its evident backing of Israel throughout the course of the 14 months?"Unfortunately, history offers little reassurance.
The \textbf{Genocide} \textbf{Convention} is not the only framework under which Washington could face scrutiny.
The Rome Statute of the \textbf{International} Criminal \textbf{Court} ( ICC ) also \textbf{recognises} complicity in \textbf{genocide} as a prosecutable offence.
In theory, the ICC \textbf{prosecutor} could bring charges against US officials.
Yet, as Lynk highlights, the decisive hurdle remains the lack of political will."What happened with South Africa ' s application against Israel breaks new ground,"Lynk observes, referencing recent \textbf{legal} efforts to hold Israel accountable for its role in the ongoing \textbf{genocidal} war in Gaza.
The ICJ is currently hearing two major \textbf{cases} linked to the Gaza war : one from South Africa accusing Israel of \textbf{genocide} through direct killings and the deprivation of essential resources, and another from Nicaragua challenging Germany ' s arms supplies to Israel.
South Africa ' s \textbf{case}, filed in late 2023 in The Hague, accuses Israel of failing to uphold its \textbf{obligations} under the 1948 \textbf{Genocide} \textbf{Convention}.
The top \textbf{court} in May ordered Israel to halt its offensive in the southern Gaza city of Rafah, marking the third time the \textbf{court} ' s 15-judge panel issued preliminary orders aimed at curbing the death toll and addressing humanitarian suffering in the besieged enclave.
The \textbf{case} has garnered \textbf{international} attention, with several other countries -- including T? rkiye, Nicaragua, Palestine, Spain, Mexico, Libya, and Colombia -- joining as parties since public hearings began in January.
In a parallel \textbf{case}, Nicaragua filed a lawsuit against Germany at the ICJ, highlighting how states that enable atrocities can also be held accountable.
Nicaragua requested \textbf{provisional} \textbf{measures} against Germany, citing its"participation in the ongoing \textbf{plausible} \textbf{genocide} and serious breaches of \textbf{international} humanitarian \textbf{law} and other peremptory norms of general \textbf{international} \textbf{law} occurring in Gaza."Adding to these developments, last week, the \textbf{International} Criminal \textbf{Court} ( ICC ) issued landmark arrest \textbf{warrants} for Israeli Prime Minister Benjamin Netanyahu and former Defence Minister Yoav Gallant.
This marked a significant shift in the enforcement of \textbf{international} \textbf{law}, signaling that no leader is beyond accountability.
Although \textbf{international} \textbf{mechanisms} have been unable to prevent \textbf{genocide} over the past year, these \textbf{legal} actions represent critical steps toward justice.
Should the ICJ rule that Israel is guilty of \textbf{genocide}, it could set a transformative precedent, potentially paving the way for \textbf{cases} against complicit states like the US, according to Lynk.
Such a move, however, would require unprecedented courage and unity from the global community."I think the whole world will be watching to see whether or not the \textbf{court}, in a couple of years ' time, makes a finding of \textbf{genocide},"Lynk reflects."And that may open further doors for action against countries where there may be either \textbf{genocide} or complicity in \textbf{genocide}."

% matched lemmas: case, convention, court, crime, genocidal, genocide, international, law, legal, measure, mechanism, obligation, plausible, prosecutor, provisional, recognise, resolution, unconditional, veto, violation, warrant
\end{textsample}
